\documentclass[11pt]{article}

% ---- Fonts & encoding ----
\usepackage[T1]{fontenc}
\usepackage[utf8]{inputenc}
\usepackage{mathptmx}

% ---- Page geometry ----
\usepackage[top=1in, bottom=1in, left=1in, right=1in]{geometry}

% ---- Packages ----
\usepackage{tikz}
\usetikzlibrary{shapes.geometric, arrows.meta, positioning, fit, backgrounds}
\usepackage{booktabs}
\usepackage{enumitem}
\usepackage{microtype}
\usepackage{setspace}
\usepackage{titlesec}
\usepackage{hyperref}
\usepackage{float}
\usepackage{indentfirst}

% ---- Section formatting ----
\titleformat{\section}{\normalfont\large\bfseries}{\thesection.}{0.5em}{}
\titleformat{\subsection}{\normalfont\normalsize\bfseries}{\thesubsection}{0.5em}{}
\titlespacing*{\section}{0pt}{10pt}{4pt}
\titlespacing*{\subsection}{0pt}{8pt}{2pt}

% ---- Spacing ----
\setlength{\parindent}{1em}
\setlength{\parskip}{3pt}

% ---- Header info ----
\title{\textbf{Offloaded Spatial Audio Rendering for XR:\\
       Latency, Power, and Quality Tradeoffs}\\[6pt]
       \normalsize CS 534 IC -- Progress Report -- Spring 2026}
\author{Abhi Ramakrishnan \quad Sethu Eapen\\
        \small University of Illinois at Urbana-Champaign}
\date{}

% -----------------------------------------------------------------------
\begin{document}
\maketitle
\thispagestyle{plain}

% -----------------------------------------------------------------------
\section{Introduction}
% -----------------------------------------------------------------------

Spatial audio is an important part of XR immersion. When sound does not
match the physical space a user is in, it is immediately noticeable.
Environment-aware spatial audio addresses this by adapting to the actual
geometry and surface materials of the room, producing a Room Impulse
Response (RIR) that captures how sound propagates in that specific space.

Running this pipeline on a headset is expensive. SAMOSA (Xu et al., UIST 2025),
a recent system for on-device scene-aware spatial audio, achieves 58\,ms
to compute an updated Room Impulse Response (RIR) on a Snapdragon XR2+ Gen 2
chip. This update latency matters because it determines how quickly the
acoustic model can adapt when a user moves to a new position in a space.
SAMOSA was only evaluated on a single chip, leaving open questions about
how it performs on other hardware.

A natural response to on-device cost is offloading computation to a nearby
server, which is the approach taken by RemoteVIO (Jiang et al., MMSys 2025)
for visual-inertial odometry in ILLIXR. For spatial audio, however, offloading
introduces a different constraint: audio-visual synchronization. Prior work
(Yeregui et al., 2024) measured audio latencies of 185\,ms over Ethernet
when fully offloading the XR rendering pipeline, exceeding the 160\,ms
threshold above which audio-video desync becomes noticeable to users. Full
offloading does not appear viable for audio, but a partial offload that keeps
lightweight steps on the headset and sends only the heavy ML inference to the
server might stay within budget.

We study the latency, power, and audio quality tradeoffs of different
offloading split points for a scene-aware spatial audio pipeline, using
a desktop workstation (Intel i5-13600K, RTX 4090) as the device proxy
and a Jetson Orin AGX as the server.

% -----------------------------------------------------------------------
\section{Related Work}
% -----------------------------------------------------------------------

\subsection{SAMOSA}

Xu et al.\ present an on-device system for scene-aware spatial
audio in XR. The pipeline runs four stages: shoebox room geometry estimation
from depth data, pixel-wise surface material classification using DeepLabv3+,
room type classification using MobileNetV2 to further refine acoustic
parameters, and RIR synthesis via the Image Source Method. It achieves
58\,ms end-to-end with under 3\% CPU usage on a Snapdragon XR2+ Gen 2.
Our pipeline follows the same structure but omits the MobileNetV2 room
classification step, focusing on the geometry and material inference stages
and how they partition across device and server. Notably, SAMOSA reports
only aggregate latency and does not break down timing per stage; our
profiling infrastructure provides this breakdown.

\subsection{RemoteVIO}

Jiang et al.\ study offloading of head tracking (VIO) in an
end-to-end ILLIXR deployment, showing that careful partitioning of a
latency-sensitive XR workload can reduce headset power without unacceptable
quality loss. We follow the same experimental structure
applied to the spatial audio pipeline instead of VIO.

\subsection{Edge Rendering for XR}

Yeregui et al.\ fully offload XR rendering and audio to an edge
server, measuring 185\,ms audio latency over Ethernet and up to 324\,ms
over WiFi. They establish 160\,ms as the perceptual threshold above which
audio-visual desync becomes noticeable. Our work asks whether a smarter
partial offload can stay within that budget where full offloading cannot.

% -----------------------------------------------------------------------
\section{Approach and Evaluation Plan}
% -----------------------------------------------------------------------

\subsection{System Architecture}

We build a scene-aware spatial audio pipeline using open-source components.
A MacBook serves as the device proxy and a Jetson Orin AGX serves as the
server, connected over a local network. The pipeline is implemented in Python,
with each stage as a standalone module. On the device side, shoebox room
geometry is estimated from a PLY mesh or depth frame and binaural rendering
is handled via Steam Audio, which requires continuous access to the current
head pose. On the server side, a Python process (\texttt{server.py}) accepts
RGB frames, runs material segmentation, computes the RIR using
\texttt{pyroomacoustics}, and returns the result over a local network socket.

We run segmentation in two modes, selectable via a flag. In the first mode,
SegFormer-b0 (Xie et al., NeurIPS 2021) pretrained on ADE20K runs inference
and its output feeds directly into RIR computation. In the second mode,
DeepLabv3+ with a MobileNetV3 backbone (torchvision pretrained) runs inference
for timing purposes only; its output is discarded and materials are set from
a fixed default, since this model is trained on PASCAL VOC and does not produce
valid indoor surface labels. The second mode lets us measure what segmentation
latency would look like with a lighter, SAMOSA-class model on our hardware
without needing to retrain on ADE20K. Audio quality metrics are only reported
for the SegFormer mode.

\subsection{Offloading Configurations}

\begin{table}[H]
\small\centering
\caption{Offloading configurations.}
\label{tab:configs}
\begin{tabular}{@{}llll@{}}
\toprule
\textbf{Config} & \textbf{On Device} & \textbf{On Server} & \textbf{Expected tradeoff}\\
\midrule
A (baseline) & Geometry + Materials + RIR & Nothing                & Highest power, lowest latency\\
B (partial)  & Geometry + RIR             & Materials              & Moderate savings\\
C (heavy)    & Geometry                   & Materials + RIR        & Large savings, latency risk\\
D (full)     & Nothing                    & Geometry + Materials + RIR & Lowest power, likely too slow\\
\bottomrule
\end{tabular}
\end{table}

Binaural rendering via Steam Audio runs on device in all configurations, as
it requires continuous access to the current head pose and the local audio
source. Config A reproduces the SAMOSA on-device baseline. Configs B and C
are partial offloads at different granularities. Config D sends all scene
processing to the server and corresponds to the full-offload case studied
by Yeregui et al., which we expect to exceed the perceptual budget.

\subsection{Metrics}

For each configuration we measure:

\begin{itemize}[noitemsep]
  \item \textbf{End-to-end latency}: time from camera frame capture to
        rendered audio output, compared against the 160\,ms threshold.
  \item \textbf{Per-step latency}: breakdown across depth estimation,
        material inference, RIR computation, and rendering.
  \item \textbf{Network bandwidth}: bytes transmitted per frame for each config.
  \item \textbf{Audio quality}: RIR similarity against a ground truth RIR
        computed by running \texttt{pyroomacoustics} directly on ScanNet's
        known room geometry, using reverberation time ($T_{60}$) and
        normalized echo density as metrics.
\end{itemize}

Experiments run on a desktop workstation (Intel i5-13600K, RTX 4090) and
Jetson Orin AGX connected over a local network, with additional WiFi runs
to test sensitivity to network conditions.

% -----------------------------------------------------------------------
\section{Methodology}
% -----------------------------------------------------------------------

The pipeline is implemented in Python with five standalone modules:
\texttt{shoebox.py} (PLY loading, PCA alignment, AABB room estimation),
\texttt{segmentation.py} (SegFormer-b0 via HuggingFace Transformers, ADE20K class mapping),
\texttt{scene\_classifier.py} (rule-based acoustic preset selection),
\texttt{acoustics.py} (pyroomacoustics ShoeBox RIR computation), and
\texttt{render.py} (FFT convolution via scipy). A \texttt{profiler.py} module
wraps every stage with \texttt{perf\_counter} timers to produce per-stage latency
breakdowns. The full pipeline is orchestrated by \texttt{pipeline.py}, which runs
shoebox estimation and material segmentation in parallel via \texttt{ThreadPoolExecutor}.

SegFormer model weights are cached at the module level so they are loaded from
disk only once per server process, reflecting the deployment scenario where the
server handles repeated RIR update requests without restarting.

Currently, the pipeline takes a pre-made PLY mesh and a static RGB render as
inputs, with source and listener positions set to fixed offsets within the room.
The next step is to replace these with inputs from ScanNet, a dataset of
RGB-D indoor scans where each scene includes per-frame depth, RGB, and camera
poses alongside a ground-truth PLY mesh. With ScanNet, the depth frame drives
shoebox estimation directly and the camera pose provides the real listener
position, removing the hardcoded placement. The ground-truth PLY mesh serves
as a reference: running pyroomacoustics on it produces a reference RIR that
we compare against our pipeline's output using $T_{60}$ and normalized echo
density as audio quality metrics.

% -----------------------------------------------------------------------
\section{Preliminary Results}
% -----------------------------------------------------------------------

\subsection{Experimental Setup}

All results reported here were collected on a desktop workstation running
Linux (CachyOS, kernel 6.19.10). The hardware used is summarized in
Table~\ref{tab:hardware}. Prior measurements taken on a MacBook are
discarded due to inconsistent thermal throttling behaviour across runs.

\begin{table}[H]
\small\centering
\caption{Evaluation hardware.}
\label{tab:hardware}
\begin{tabular}{@{}ll@{}}
\toprule
\textbf{Component} & \textbf{Specification}\\
\midrule
CPU  & Intel Core i5-13600K (13th Gen)\\
GPU  & NVIDIA GeForce RTX 4090 (24\,GB VRAM)\\
RAM  & 62\,GB DDR5\\
OS   & Linux (CachyOS), kernel 6.19.10\\
Runtime & Python 3.12, PyTorch 2.x, CUDA\\
\bottomrule
\end{tabular}
\end{table}

\subsection{Config A: SegFormer Segmentation Mode}

Table~\ref{tab:results} shows per-stage latency for Config A (all stages
on one device) using SegFormer-b0 for material segmentation. All values are
warm-run medians over 10 consecutive pipeline executions with models and
weights pre-loaded in GPU memory.

\begin{table}[H]
\small\centering
\caption{Per-stage latency, Config A, SegFormer mode (warm runs, RTX 4090).}
\label{tab:results}
\begin{tabular}{@{}lrr@{}}
\toprule
\textbf{Stage} & \textbf{Range (ms)} & \textbf{Median (ms)}\\
\midrule
Shoebox estimation     & 5\,--\,9    &  7.2\\
Material segmentation  & 89\,--\,110 & 100.2\\
Scene classification   & $<$1        & $<$1\\
RIR synthesis (ISM)    & 10\,--\,29  & 15.8\\
Audio rendering        & 1\,--\,4    &  1.2\\
\midrule
\textbf{Total (parallel phase 1)} & & \textbf{$\approx$117}\\
\bottomrule
\end{tabular}
\end{table}

Material segmentation dominates at 100\,ms, accounting for approximately
86\% of total pipeline time. The total of $\approx$117\,ms exceeds
SAMOSA's 58\,ms end-to-end figure. This gap is expected: SAMOSA uses a
custom-trained 9-class MobileNetV2 optimised for on-device inference,
whereas SegFormer-b0 is a transformer-based model trained on 150 ADE20K
classes and not purpose-built for latency. Shoebox estimation (7\,ms) and
RIR synthesis (16\,ms median) are both well below the SAMOSA equivalents.

\subsection{Config A: SAMOSA Emulation Mode}

Because SAMOSA's trained segmentation model is not publicly available
(the paper is a closed-source Google Research publication), we implement a
SAMOSA emulation mode to produce a timing-comparable baseline. In this mode,
DeepLabv3+ with a MobileNetV3-Large backbone (torchvision pretrained, COCO)
replaces SegFormer during the segmentation slot. Its output is discarded;
material labels are loaded from a one-time SegFormer cache computed
offline via \texttt{--init-segmentation}. This approximates SAMOSA's
inference cost (same dense-prediction architecture family, comparable parameter
count) without requiring access to their proprietary training data. Audio
quality metrics are only reported for SegFormer mode.

\begin{table}[H]
\small\centering
\caption{Per-stage latency, Config A, SAMOSA emulation mode (warm runs, RTX 4090).}
\label{tab:results-samosa}
\begin{tabular}{@{}lrr@{}}
\toprule
\textbf{Stage} & \textbf{Range (ms)} & \textbf{Median (ms)}\\
\midrule
Shoebox estimation           & 5\,--\,9   & 7.2\\
Segmentation (DeepLabv3+MNv3) & 8\,--\,8  & 7.8\\
Scene classification         & $<$1       & $<$1\\
RIR synthesis (ISM)          & 10\,--\,29 & 15.8\\
Audio rendering              & 1\,--\,4   & 1.2\\
\midrule
\textbf{Total (parallel phase 1)} & & \textbf{$\approx$25}\\
\bottomrule
\end{tabular}
\end{table}

In SAMOSA emulation mode, segmentation drops to 7.8\,ms, bringing the
total pipeline latency to approximately 25\,ms --- well under both the
SAMOSA 58\,ms reference and the 160\,ms perceptual budget. This confirms
that the pipeline structure is sound and that segmentation is the dominant
cost. Table~\ref{tab:comparison} compares the two modes against SAMOSA.

\begin{table}[H]
\small\centering
\caption{Pipeline latency comparison (warm runs).}
\label{tab:comparison}
\begin{tabular}{@{}lrrr@{}}
\toprule
\textbf{System} & \textbf{Segmentation} & \textbf{RIR} & \textbf{Total}\\
\midrule
SAMOSA (Xu et al.)         & $\approx$10\,ms & $\approx$30\,ms & 58\,ms\\
Ours (SegFormer mode)      & 100\,ms          & 16\,ms          & $\approx$117\,ms\\
Ours (SAMOSA emul.\ mode)  & 7.8\,ms          & 16\,ms          & $\approx$25\,ms\\
\bottomrule
\end{tabular}
\end{table}

The SAMOSA total (58\,ms) includes ISM early reflections plus a spectral
late-reverberation tail; our RIR stage uses ISM only, which accounts for
the 16\,ms vs.\ $\approx$30\,ms difference. The segmentation bottleneck
directly motivates the offloading design: Config~B (offload segmentation
to a GPU server) is expected to reduce device-side cost to under 10\,ms,
well within budget even after accounting for network round-trip time.

% -----------------------------------------------------------------------
\section{Changes from Original Proposal}
% -----------------------------------------------------------------------

\textbf{Testbed: MacBook + Jetson Orin without ILLIXR.}
The original proposal targeted ILLIXR as the device-side testbed. On instructor
feedback, full ILLIXR integration was deemed overly ambitious for the available
timeline. We instead run the pipeline directly on a MacBook (device proxy) and
Jetson Orin AGX (server) without the ILLIXR layer. This keeps the hardware
setup close to the proposal while removing the integration overhead. Power
measurements via \texttt{tegrastats} on the Jetson remain feasible with this setup.

\textbf{Segmentation model: SegFormer-b0 instead of DeepLabv3+.}
The proposal specified DeepLabv3+ following SAMOSA. SAMOSA's model was
custom-trained on COCO-Stuff and is not publicly available, and no pretrained
DeepLabv3+ trained on ADE20K exists off the shelf. We use SegFormer-b0
pretrained on ADE20K, which provides equivalent pixel-level indoor surface
classification (150 classes) with a comparable parameter count (3.75M).
Since the segmentation model is offloaded to the server in Configs B--D,
its specific architecture does not affect the offloading comparison.

% -----------------------------------------------------------------------
\section{Remaining Timeline}
% -----------------------------------------------------------------------

\begin{table}[H]
\small\centering
\caption{Remaining project milestones.}
\label{tab:timeline}
\begin{tabular}{@{}ll@{}}
\toprule
\textbf{Task} & \textbf{Target}\\
\midrule
\texttt{server.py} socket layer; device/server split working end-to-end & Apr 10\\
Config B implemented and measured (offload segmentation)                & Apr 13\\
Configs C and D implemented and measured                                & Apr 16\\
ScanNet ground truth RIR comparison; audio quality metrics              & Apr 19\\
Full experiment runs across all configs; latency and bandwidth tables   & Apr 21\\
Results analysis; bottleneck identification; plots and tables           & Apr 24\\
Final report and presentation                                           & Apr 28\\
\bottomrule
\end{tabular}
\end{table}

\end{document}
